\section{Discussion}
\label{sec:discussion}

\para{What do the results mean?} The main result is not that dilution makes hidden structure invisible. It is that dilution breaks exact recovery before it breaks detection. That distinction matters for any monitor that relies on a prompted LLM alone. A model may correctly suspect that a passage is unusual, yet fail to reconstruct the payload needed for attribution, filtering, or automated response.

\para{Why is \gptfourone stronger here?} On this benchmark, \gptfourone is more reliable than \gptfive for extraction, especially on \densekw\ and \dilutedkw. We do not treat this as a broad ranking claim between the models. The result is prompt- and API-specific, and it likely depends on how each model balances reasoning behavior with constrained JSON output. Still, under the exact protocol used here, \gptfourone preserves more payload information once the signal becomes sparse.

\para{What failures dominate?} The processed outputs show three recurring errors. First, diluted passages often trigger false negatives in extraction, with the model returning \texttt{NONE} despite detecting some suspicious structure. Second, near-miss decodes are common, including one-letter substitutions or transpositions such as \texttt{PRAIRA}, \texttt{QUARRY}, \texttt{EMBERB}, and \texttt{NELBUL}. Third, longer diluted passages become brittle when distractor forms compete with the true trigger. These are not random errors. They are exactly the kinds of failures one would expect when a model tracks local clues but loses the global decoding path.

\para{Limitations.} This study is deliberately small and controlled. Each scheme-model comparison uses only 12 positive examples, so uncertainty remains high. The benchmark is synthetic, even though the cover text comes from real corpora, and our negatives are clean matched controls rather than a broad sample of natural text. The model is told the candidate hiding schemes, so this is not open-ended steganalysis. Finally, the \gptfive validation fix shows that API behavior can materially affect results; our conclusions should therefore be read as claims about this benchmark, prompt set, and API configuration rather than universal model properties.

\para{Broader implications.} These findings sharpen the current safety picture from recent frontier-model studies \cite{karpov2025potentials,zolkowski2025early}. A sparse covert channel does not need to be fully undetectable to be operationally problematic. It may be enough that a monitor can flag suspicion but cannot reliably decode the message. That suggests future monitoring systems should combine prompted models with stronger search, rule induction, or programmatic decoding components rather than expecting a single LLM call to do both detection and extraction robustly.
